\section{SWiRL: Synthetic Data y Multi-Step RL}

\subsection{Objetivos de diseño}

Las tareas de múltiples pasos con herramientas requieren que el modelo aprenda: cuándo razonar, cuándo invocar una herramienta, cómo consumir la observation, cómo recuperarse después de un error, y cuándo detenerse y dar la respuesta final.

\videofigure{figures/swirl-motivation.jpg}{Las tareas reales suelen requerir razonamiento de múltiples pasos y uso de herramientas.}{00:50:30--00:53:02}

\videofigure{figures/swirl-goals.jpg}{SWiRL busca aprender la selección de herramientas, el consumo de retroalimentación, la recuperación de errores y el momento de detenerse.}{00:53:02--00:54:34}

La restricción clave de SWiRL es: al entrenar con RL, en la medida de lo posible no invocar herramientas reales en línea, porque las herramientas son lentas, fallan con facilidad y son difíciles de paralelizar. Primero genera fuera de línea trayectorias completas con observation, y después realiza sobre esos datos fijos un RL a nivel de paso (step-wise).

\subsection{Etapa 1: síntesis de trayectorias de múltiples pasos fuera de línea}

Para cada prompt, el teacher model genera una sola action a la vez: puede ser reasoning, una llamada a herramienta o la respuesta final. Si se invoca una herramienta, el entorno fuera de línea la ejecuta y escribe la observation de vuelta en el contexto, y entonces continúa al siguiente paso. Un LLM judge asigna un process score a cada paso.

\videofigure{figures/swirl-data-generation.jpg}{SWiRL genera y filtra fuera de línea trayectorias de reasoning--tool de múltiples pasos.}{00:54:34--00:58:32}

Dos formas principales de filtrado:

\begin{itemize}
    \item \textbf{Process filtered}: se conservan las trayectorias en las que cada paso es considerado razonable por el judge;
    \item \textbf{Outcome filtered}: solo se conservan las trayectorias cuya respuesta final es correcta, sin importar la puntuación de los pasos intermedios;
    \item también se puede exigir que pasen tanto el criterio de process como el de outcome.
\end{itemize}

\subsection{Etapa 2: Step-Wise Reinforcement Learning}

Durante el entrenamiento ya no se invocan herramientas, sino que se reproduce la tool response de los datos fuera de línea. El modelo elige una action en cada estado, y la policy se actualiza usando la reward de paso.

\videofigure{figures/swirl-step-rl.jpg}{SWiRL ejecuta aprendizaje por refuerzo a nivel de paso sobre las trayectorias fuera de línea.}{00:58:32--01:03:54}

\videofigure{figures/swirl-objective.jpg}{El objetivo de RL de múltiples pasos de SWiRL acumula la reward para cada par estado--acción.}{01:03:54--01:05:16}

Se puede resumir como:

$$
J(\theta)=\mathbb E_{(s_t,a_t)\sim\mathcal D}
\left[w_t\log\pi_\theta(a_t\mid s_t)\right],
$$

donde $w_t$ está determinado por la reward de paso o la advantage. Las action de alta calidad ven aumentada su probabilidad, y las de baja calidad se ven suprimidas.

\begin{knowledgebox}{Ventaja de ingeniería de las trayectorias de herramientas fuera de línea}
La ejecución de herramientas y la anotación del judge pueden hacerse en paralelo por adelantado; el entrenamiento de RL solo lee un contexto fijo, con lo que se obtiene mayor rendimiento, un entorno más estable y experimentos más reproducibles. Pero el modelo no puede explorar acciones nuevas que no estén en el conjunto de datos, y su desempeño está limitado por la cobertura fuera de línea.
\end{knowledgebox}

\subsection{Etapa 3: razonamiento real de múltiples pasos}

En el despliegue, el modelo vuelve a conectarse con herramientas reales y, siguiendo la política obtenida en el entrenamiento, itera entre reasoning, tool use y la respuesta final.

\videofigure{figures/swirl-inference.jpg}{En la etapa de inferencia se vuelven a habilitar las herramientas reales y se ejecuta un ciclo cerrado de múltiples pasos.}{01:05:16--01:06:16}

\subsection{Por qué el Process Filtering resulta mejor}

En los experimentos, los datos filtrados únicamente por calidad de process superaron a los datos que solo conservan la respuesta final correcta, o que exigen simultáneamente que el proceso y el resultado sean correctos.

\videofigure{figures/swirl-filtering.jpg}{Las trayectorias sintéticas filtradas por process (process-filtered) producen la mayor ganancia de entrenamiento.}{01:06:16--01:08:18}

La razón es que el outcome filtering solo conserva los problemas que el teacher ya sabe resolver, lo que fácilmente forma una «distribución estrecha de trayectorias exitosas»; el process filtering, en cambio, también conserva problemas difíciles en los que los pasos locales son correctos pero el resultado final falla, y así ofrece al modelo nuevos componentes de razonamiento y señales de recuperación de errores.

\begin{importantbox}{La corrección local puede aportar más cobertura de aprendizaje que el éxito final}
Una trayectoria que al final falla puede seguir conteniendo una buena descomposición del problema, invocaciones correctas de herramientas o evidencia intermedia útil. Si solo se descarta la trayectoria completa según la reward final, esas habilidades reutilizables también se eliminan. El process filtering cambia la selección de datos de «éxito o fracaso del problema completo» a «valor de cada paso».
\end{importantbox}

\subsection{Generalización entre tareas y entre herramientas}

Al entrenar con GSM8K + calculator, el modelo también mejora en HotpotQA + search; y viceversa. Esto indica que el modelo no solo memoriza una API específica, sino que aprende una política general de gestión de estado en múltiples pasos y de «cuándo recurrir a una herramienta».

\videofigure{figures/swirl-generalization.jpg}{SWiRL muestra generalización entre herramientas entre las tareas de matemáticas y de preguntas y respuestas de múltiples saltos.}{01:08:18--01:10:02}

\videofigure{figures/swirl-data-scaling.jpg}{Aumentar la escala de los datos de entrenamiento sintéticos sigue mejorando el desempeño tanto dentro como fuera de la distribución.}{01:10:02--01:11:34}

\videofigure{figures/process-correctness.jpg}{La mejora del process reward fuera de la distribución indica que la ganancia proviene de un mejor proceso de múltiples pasos.}{01:11:34--01:12:20}

\subsection{Limitaciones}

\begin{itemize}
    \item El sesgo del LLM judge se transmite al process reward;
    \item la observation fuera de línea impide que la policy experimente errores nuevos de las herramientas;
    \item el costo de generar trayectorias es alto y la cobertura puede ser insuficiente;
    \item cuando el schema de una herramienta nueva difiere demasiado del de las herramientas de entrenamiento, la generalización sigue fallando.
\end{itemize}

\subsection{Resumen de la sección}

SWiRL convierte la interacción en línea con herramientas en synthetic trajectory fuera de línea, y después aprende una política de múltiples pasos con RL a nivel de paso. El resultado más importante es que filtrar por calidad de proceso favorece más la generalización que exigir únicamente la corrección del resultado final, lo que indica que los datos de entrenamiento deben conservar las capacidades locales reutilizables.
